\subsection{Spectral characterization of Riesz endomorphisms}

Let E be a complex Banach space and u an endomorphism of E. Recall that we denote by \(\mathrm{Sp}(u)\) the spectrum of u relative to the unital Banach algebra \(\mathcal{L}(\mathrm{E})\) (cf. § 7 of I, p. 127).

Suppose that 0 is an isolated point of \(\mathrm{Sp}(u)\) ; recall that we then denote by \(e_0(u)\) the spectral projector associated with \(u\) and with the open and closed subset \(\{0\}\) of the spectrum of \(u\) (cf. n°3 of I, p. 131). We have \(e_0(u) = f(u)\) for every germ of holomorphic function \(f\) in the neighborhood of \(\mathrm{Sp}(u)\) which equals 1 in a neighborhood of 0 and vanishes in a neighborhood of \(\mathrm{Sp}(u) \setminus \{0\}\) .

By passing to subspaces, the endomorphism \(u\) induces a quasi-nilpotent endomorphism of the image of \(e_0(u)\) , whose spectrum is reduced to \(\{0\}\) , and an automorphism of the kernel of \(e_0(u)\) , whose spectrum is \(\mathrm{Sp}(u) \setminus \{0\}\) (loc. cit.).

PROPOSITION 14. --- Let E be a complex Banach space. For an element u of \(\mathcal{L}(\mathrm{E})\) to be a Riesz endomorphism of E, it is necessary and sufficient that one of the following two mutually exclusive conditions be satisfied:

\begin{enumerate}
\def\labelenumi{(\roman{enumi})}
\item
  The endomorphism u of E is an automorphism of E;
\item
  The point 0 is an isolated point of \(\operatorname{Sp}(u)\) and the projector \(e_0(u)\) has finite rank.
\end{enumerate}

When condition (ii) is satisfied, the image of \(e_{0}(u)\) is the nilespace of u and the kernel of \(e_{0}(u)\) is the conilespace of u.

Every automorphism of E is a Riesz endomorphism of E (III, p. 47, remark 1). If u satisfies condition (ii), the kernel F of \(e_{0}(u)\) is a closed vector subspace of finite codimension of E, stable under u. Let \(u_{F}\) be the endomorphism of F deduced from u. Its spectrum is contained in \(\mathbf{C}\setminus\{0\}\) (n°3 of I, p. 131), hence \(u_{F}\) is an automorphism of F. It follows that u is a Riesz endomorphism of E (III, p. 48, prop. 8).

Conversely, let \(u\) be a Riesz endomorphism of E. Denote by N its nilespace and by I its conilespace. By prop. 6 of III, p. 47, the space E is the topological direct sum of N and I, and the map \(u\) defines by passage to subspaces a nilpotent endomorphism \(u_{\mathrm{N}}\) of N and an automorphism \(u_{\mathrm{I}}\) of I. In particular, we have \(\mathrm{Sp}(u_{\mathrm{N}}) \subset \{0\}\) and \(0 \not\in \mathrm{Sp}(u_{\mathrm{I}})\) . If N is zero, then I = E and \(u\) is an automorphism of E. Otherwise, 0 is an isolated point of \(\mathrm{Sp}(u)\) and \(e_0(u)\) is the projector with image N and kernel I (prop. 2 of I, p. 129); in particular, it is of finite rank.

Remark. --- Let E be a real Banach space and let u be an endomorphism of E. Suppose that 0 is an isolated point of the complex spectrum of u, that is, of the spectrum of the endomorphism \(1 \otimes u\) of the complete normable complexified space \(\mathrm{E}_{(\mathbf{C})}\) of E (I, p. 85, n° 13). The subset \(\{0\}\) of the complex spectrum of \(u\) is open and closed, and invariant under conjugation; denote by \(e_0(u) \in \mathcal{L}(\mathrm{E})\) the spectral projector associated with it (n° 13 of I, p. 85). Proposition 14 holds mutatis mutandis in this setting.

